\documentclass[10pt,a4paper]{amsart}
\usepackage[margin=19mm]{geometry}
\usepackage{amsmath,amssymb,mathtools}
\usepackage[scr=rsfs,cal=euler]{mathalfa}
\usepackage{xfrac}
\usepackage[unicode,hidelinks,bookmarks=false]{hyperref}
\newcommand{\ie}{i.e.\ }
\newcommand{\cf}{cf.\ }
\newcommand{\resp}{respectively\ }
\newcommand{\Subsection}{\subsection}
\providecommand{\nobreakdash}{\nobreak\hbox{-}}
\setlength{\emergencystretch}{2em}
\multlinegap0pt
\usepackage{ulem}
\newtheorem{theorem}{Theorem}
\newtheorem{lemma}[theorem]{Lemma}
\newtheorem{proposition}[theorem]{Proposition}
\newtheorem{corollary}[theorem]{Corollary}
\newtheorem{conjecture}[theorem]{Conjecture}
\newtheorem{claim}[theorem]{Claim}
\newtheorem{question}[theorem]{Question}
\newtheorem{remark}[theorem]{Remark}
\newtheorem{example}{Example}
\newcommand{\bExample}{\begin{example}}
\newcommand{\eExample}{\end{example}}
\newcommand{\bEnum}{\begin{enumerate}}
\newcommand{\eEnum}{\end{enumerate}}
\newcommand{\bItemize}{\begin{itemize}}
\newcommand{\eItemize}{\end{itemize}}
\newcommand{\N}{\mathbb N}
\newcommand{\Z}{\mathbb Z}
\newcommand{\R}{\mathbb R}
\newcommand{\bd}{\mathrm{bd}} %boundary
\newcommand{\sm}{\setminus}
\newcommand{\st}{: \;}
\newcommand{\ol}{\overline}
\newcommand{\ul}{\underline}
\newcommand{\inv}{^{-1}}
\newcommand{\ints}[1]{\{0,\cdots,#1\}}
\newcommand{\ty}{\nabla\mathrm{Y}}
\newcommand{\yt}{\mathrm{Y}\nabla}
\newcommand{\PP}{\mathcal P}
\newcommand{\dsp}{\displaystyle}
\newcommand{\arb}{\addtocounter{arb}{1}{\arabic{arb}}}
\newcommand{\barb}{\addtocounter{arb}{1}{\bf\arabic{arb}}}
\newcommand{\bibtitle}[1]{{#1}}
\begin{document}
\title[Three thousand obstructions to knotless embedding]{Three thousand obstructions to knotless embedding}
\author{CSU, Chico; Thomas W Mattman; Andrei Pavelescu}
\date{}
\maketitle
\noindent\emph{Source and licence.} Three thousand obstructions to knotless embedding. Version arXiv:2607.01414v1 (CC BY 4.0). This material is distributed under the Creative Commons Attribution 4.0 International licence, http://creativecommons.org/licenses/by/4.0/, as recorded in the arXiv OAI licence metadata for this exact version and shown on the abstract page https://arxiv.org/abs/2607.01414v1. Full rights notice: licenses/arxiv-2607.01414-CC-BY-4.0.txt.\par\smallskip
\noindent\textit{Editorial scope.} Counted unit: the original \texttt{proposition} environment \texttt{prop-mmik3cut} at source lines 407--413 together with its complete proof at lines 415--465; the delivered document keeps source lines 338--466 so that every referenced label and every citation resolves inside it. Native editable source is the paper's own concatenated TeX; the AMS wrapper is editorial formatting only. No publisher class, upstream script or unrelated artwork is redistributed.
\section{Claws in MMIK graphs}
\label{sec:claws}

%\begin{lemma}
%\label{lem-mmikclaw}
%Vertices of degree 3 in MMIK graphs
%are not in 3-cycles.
%\end{lemma}

%\begin{proof}
%Let $x$ be a vertex of degree 3 in a MMIK graph $G$.
%Let $a, b, c$ denote the neighbors of $x$.
%For contradiction, suppose $x$ is in a 3-cycle.
%Then, without loss of generality, $ab$ is an edge of $G$.
%Since $G$ is MMIK, $G-ab$ has a knotless embedding,  $\eta$, in $\R^3$.
%(We denote $\eta(G-ab)$ as just $\eta$.)
%Add $ab$ to $\eta$ by embedding it 
%``close and parallel'' to the path $axb$.
%Call this new embedding $\nu$.
%We will show $\nu$ is knotless.

%Let $C$ be a cycle in $\nu$.
%If $ab$ is not in $C$, then $C$ is in $\eta$ and hence a trivial knot.
%So assume $ab$ is in $C$.
%If $x \not \in C$,
%then $C$ is isotopic to the cycle $C -ab + ax + xb \s%ubset \eta$,
%which is a trivial knot.
%So assume $x \in C$.
%If $C = axba$, then $C$ is a trivial knot.
%If $C \ne axba$, then, without loss of generality,
%$C$ contains the path $abxc$.
%So $C$ is isotopic to the cycle
%$C -ab -bx + ax \subset \eta$,
%which is a trivial knot.
%\end{proof}

In this section we discuss the structure of a MMIK graph near a vertex of degree three or four as well as the connectivity of these graphs.
We begin by observing that the neighborhood of a degree three vertex induces a $K_{3,1}$ or claw subgraph.
The following result refers to properties, such as IK and non-planar,
that are preserved by the $\ty$ move and by contracting an edge on a degree two vertex.

\begin{lemma}
\label{lem-mmikclaw}
Let $\PP$ be any graph property that is preserved by the $\ty$ move
and by contracting an edge on a degree two vertex.
Then in every graph that is minor minimal with respect to property P,
every vertex of degree 3 has a claw neighborhood.
\end{lemma}
\begin{proof}
Suppose $G$ has property $\PP$,
$x$ has degree 3 in $G$,
its neighbors are $a, b, c$,
and $ab$ is an edge of $G$.
Do a $\ty$ move on $xab$.
Let $y$ denote the center of the new Y.
Now contract $xy$.
The graph we get is isomorphic to $G - ab$
and has property $\PP$.
So $G$ is not minor minimal with respect to P.
\end{proof}

\begin{corollary}
\label{cor-mmikclaw}
Vertices of degree 3 in MMIK graphs
are not in 3-cycles.
\end{corollary}

Equivalently, since MMIK graphs have minimum degree three, every vertex in a 3-cycle has degree at least four.


\begin{proposition}
\label{prop-mmik3cut}
If $G$ is a MMIK graph and $\{a, b, c\}$ is a 3-cut in $G$,
then $G$ has a vertex $x$ such that
the subgraph of $G$ induced by $\{x,a,b,c\}$ consists of
the edges $xa, xb, xc$ only.
\end{proposition}

\begin{proof}
Let $G_1, G_2$ be distinct connected components of $G -a-b-c$.
If, for some $i \in \{1,2\}$, the graph $G_i$ has only one vertex,
then we are done  by Lemma~\ref{lem-mmikclaw}.
So assume each $G_i$ has at least two vertices; we will show $G$ is nIK.
Let $H_i$ be the proper minor of $G$ obtained by contracting $G_i$ to one vertex, $x_i$.
Since $G$ is MMIK, $H_i$ has a knotless embedding, $\eta_i$, in $\R^3$.
Isotope $\eta_1$ so that it lies in $\R^3_+ := \{(x,y,z) \st z \ge 0\}$
and intersects the $xy$-plane in precisely $\{x_1a, x_1b, x_1c\}$.
Similarly, isotope $\eta_2$ to lie in $\R^3_- := \{(x,y,z) \st z \le 0\}$,
but with the added condition that $x_2 a, x_2 b, x_2 c$ coincide, respectively, 
with   $x_1 a, x_1 b, x_1 c$.
Let $\eta = \eta_1 \cup \eta_2$,
with $x = x_1 = x_2$.
Then $\eta - x =G$.
So, to show $G$ is nIK, it suffices to show $\eta$ is knotless.

Let $C$ be a cycle in $\eta$.
If $C \subset \eta_i$ for some $i$,
then $C$ is a trivial knot since $\eta_i$ is knotless.
So assume $C$ has 
at least one vertex above, and at least one vertex below, the $xy$-plane.

Case 1. 
$x \not \in C$.
Then we can assume 
$C= (a, v_1, \cdots, v_s, b, w_1, \cdots, w_t, a)$
for some $s \ge 1$ and $t \ge 1$, 
with $v_1, \cdots, v_s$ above the $xy$-plane
and $w_1, \cdots, w_t$ below the $xy$-plane.
The cycles 
$(a, v_1, \cdots, v_s, b, x, a)$ and
$(b, w_1, \cdots, w_t, a, x, b)$
are trivial knots since they lie in $\eta_1$ and $\eta_2$, respectively.
These two cycles intersect in the path $axb$,
so their connected sum, which is $C$, is a trivial knot.

Case 2. 
$x \in C$.
Then we can assume 
$C= (a, v_1, \cdots, v_s, b, x, c, w_1, \cdots, w_t, a)$
for some $s \ge 1$ and $t \ge 1$, 
with $v_1, \cdots, v_s$ above the $xy$-plane
and $w_1, \cdots, w_t$ below the $xy$-plane.
The cycles 
$(a, v_1, \cdots, v_s, b, x, a)$ and
$(c, w_1, \cdots, w_t, a, x, c)$
are trivial knots since they lie in $\eta_1$ and $\eta_2$, respectively.
These two cycles intersect in the edge $ax$,
so their connected sum, which is $C$, is a trivial knot.
\end{proof}


\end{document}
